\documentclass[conference]{IEEEtran}
\IEEEoverridecommandlockouts

\usepackage{cite}
\usepackage{amsmath,amssymb,amsfonts}
\usepackage{algorithmic}
\usepackage{graphicx}
\usepackage{textcomp}
\usepackage{xcolor}
\usepackage{booktabs}
\usepackage{microtype}
\usepackage{url}

\def\BibTeX{{\rm B\kern-.05em{\sc i\kern-.025em b}\kern-.08em
    T\kern-.1667em\lower.7ex\hbox{E}\kern-.125emX}}

\begin{document}

\title{Generative Graph World Models for Predictive Intrusion Detection and Kill-Chain Forecasting}

\author{\IEEEauthorblockN{Raghuraj P. Singh}
\IEEEauthorblockA{\textit{Department of Artificial Intelligence \& Machine Learning} \\
\textit{B.Tech AIML}\\
raghuraj2008.cu@gmail.com}
}

\maketitle

\begin{abstract}
Modern Network Intrusion Detection Systems (NIDS) and Security Information and Event Management (SIEM) architectures rely on post-facto signature and anomaly classification on isolated flow instances ($y = f(x_t)$), resulting in zero predictive lead time ($0\text{s}$) and vulnerability to slow-and-low multi-stage attacks. This paper presents \textbf{Vanguard-GWM}, a generative Graph World Model architecture that reframes cyber defense as continuous temporal environment simulation. By slicing raw packet and flow telemetry into discrete temporal graph snapshots ($G_t$), spatial topology is compressed into low-dimensional latent state representations ($z_t$) via a Graph Attention Network (GATv2) regularized with InfoNCE contrastive self-supervision. A Causal Temporal Transformer dynamics model learns the transition distribution $P(z_{t+1} \mid z_t)$ to execute autoregressive $K$-step forward rollouts ($[\hat{z}_{t+1}, \dots, \hat{z}_{t+K}]$). Evaluated on benchmark temporal attack datasets, Vanguard-GWM achieves an $F_1$-score of \textbf{0.955}, maintains a low \textbf{1.1\% False Positive Rate}, and provides a \textbf{$+180\text{s}$ to $+600\text{s}$ predictive lead time} across standardized MITRE ATT\&CK kill-chain phases while enabling zero-risk ``What-If'' defense policy simulations.
\end{abstract}

\begin{IEEEkeywords}
Network Security, Graph Neural Networks, World Models, Attack Forecasting, MITRE ATT\&CK, Explainable AI.
\end{IEEEkeywords}

\section{Introduction}
Modern enterprise networks and Critical Information Infrastructures (CII) operate under persistent threat from sophisticated, multi-stage cyber campaigns. Advanced Persistent Threats (APTs) rarely execute monolithic exploits; instead, they progress systematically along structured attack life cycles---spanning reconnaissance, initial access, internal credential harvesting, lateral movement, and data exfiltration.

Despite substantial investments in Security Operations Centers (SOCs), existing defense frameworks suffer from two structural limitations:
\begin{enumerate}
    \item \textbf{Isolated Static Instance Triage:} Mainstream machine learning and deep learning NIDS model traffic flows independently ($y = f(x_t)$), missing the topological correlations that characterize low-and-slow lateral movement.
    \item \textbf{Post-Facto Reactive Stance:} Alerts trigger only after an anomalous flow is transmitted or a payload detonates, giving defenders zero predictive lead time ($0\text{s}$).
\end{enumerate}

To overcome these structural bottlenecks, this paper introduces \textbf{Vanguard-GWM}, a generative Graph World Model architecture that reframes cyber defense as continuous temporal environment simulation. Vanguard-GWM models the underlying network transition dynamics $P(z_{t+1} \mid z_t)$, executing autoregressive $K$-step forward rollouts to forecast evolving attack trajectories \textbf{3 to 10 minutes before execution}.

\section{Related Work}
\subsection{Machine Learning in NIDS}
Shallow classifiers (Random Forest, SVM) and deep learning models (CNN-LSTMs) evaluate isolated flows or sequential host logs. However, these models cannot capture distributed enterprise-wide topologies where attacks span multiple asynchronous hosts simultaneously.

\subsection{Graph Neural Networks for Cyber Defense}
Graph Convolutional Networks (GCNs) and Graph Attention Networks (GATs) have demonstrated success in identifying lateral movement in provenance graphs. Nevertheless, existing GNN frameworks operate in \textit{discriminative detection modes}---classifying existing graphs rather than forecasting future topological transitions.

\subsection{Model-Based World Models}
World Models decouple spatial perception from temporal transition dynamics to simulate future environments ($P(S_{t+1} \mid S_t, A_t)$). Vanguard-GWM brings this paradigm to cybersecurity by replacing pixel observation spaces with dynamic network graphs ($G_t$).

\section{System Architecture \& Formulation}

\subsection{Dual-Level Telemetry Slicing ($G_t$)}
Network traffic is partitioned into non-overlapping temporal windows $\Delta t \in [1\text{s}, 5\text{s}]$. For each time window $t$, an attributed directed graph $G_t = (V_t, E_t, \mathbf{X}_t, \mathbf{E}_t)$ is constructed where nodes $V_t$ represent active IP endpoints and edges $E_t$ represent communication sessions. Edge attributes $\mathbf{e}_{uv}$ fuse macro-flow and micro-packet telemetry:
\begin{equation}
\mathbf{e}_{uv} = \left[ \mathbf{b}_{\text{flags}} \,\|\, \log(\text{Bytes}_{uv} + 1) \,\|\, \mu(\Delta \tau_{uv}) \,\|\, \sigma^2(\Delta \tau_{uv}) \,\|\, \sigma^2(\text{TTL}_{uv}) \right]
\end{equation}
where $\Delta \tau_{uv}$ denotes packet Inter-Arrival Time (IAT) and $\sigma^2(\text{TTL}_{uv})$ captures Time-to-Live variance.

\subsection{Spatial Perception via Contrastive GATv2}
A multi-head GATv2 encoder computes dynamic edge-conditioned attention weights:
\begin{equation}
\alpha_{ij} = \frac{\exp\left(\text{LeakyReLU}\left(\mathbf{a}^\top \left[ \mathbf{W}_s \mathbf{h}_i \,\|\, \mathbf{W}_t \mathbf{h}_j \,\|\, \mathbf{W}_e \mathbf{e}_{ij} \right]\right)\right)}{\sum_{k \in \mathcal{N}(i)} \exp\left(\text{LeakyReLU}\left(\mathbf{a}^\top \left[ \mathbf{W}_s \mathbf{h}_i \,\|\, \mathbf{W}_t \mathbf{h}_k \,\|\, \mathbf{W}_e \mathbf{e}_{ik} \right]\right)\right)}
\end{equation}
Graph-level pooling synthesizes a latent state vector $\mathbf{z}_t \in \mathbb{R}^d$, regularized via self-supervised InfoNCE loss:
\begin{equation}
\mathcal{L}_{\text{InfoNCE}} = -\log \frac{\exp\left(\frac{\mathbf{z}_t^\top \mathbf{z}_t^+}{\|\mathbf{z}_t\| \|\mathbf{z}_t^+\| \cdot \tau}\right)}{\exp\left(\frac{\mathbf{z}_t^\top \mathbf{z}_t^+}{\|\mathbf{z}_t\| \|\mathbf{z}_t^+\| \cdot \tau}\right) + \sum_{j=1}^B \exp\left(\frac{\mathbf{z}_t^\top \mathbf{z}_j^-}{\|\mathbf{z}_t\| \|\mathbf{z}_j^-\| \cdot \tau}\right)}
\end{equation}

\subsection{Dynamics Engine and Autoregressive Rollout}
A Causal Temporal Transformer processes historical latent states $\mathbf{Z}_{t-w:t}$ to predict one-step transitions $\hat{\mathbf{z}}_{t+1}$. Forward $K$-step autoregressive rollouts are executed recurrently:
\begin{equation}
\hat{\mathbf{z}}_{t+k} = \text{DynamicsEngine}\left([\mathbf{z}_{t-w+k:t}, \hat{\mathbf{z}}_{t+1}, \dots, \hat{\mathbf{z}}_{t+k-1}]\right)
\end{equation}

Multi-task heads evaluate each simulated vector $\hat{\mathbf{z}}_{t+k}$ to produce:
\begin{enumerate}
    \item \textbf{Infiltration Risk:} $\text{Risk}(t+k) = \sigma(\mathbf{W}_r \hat{\mathbf{z}}_{t+k} + b_r) \in [0, 1]$
    \item \textbf{MITRE ATT\&CK Stage:} $P(\text{Stage}_{t+k} = c) = \text{Softmax}(\mathbf{W}_c \hat{\mathbf{z}}_{t+k} + \mathbf{b}_c)$
    \item \textbf{Time-to-Compromise:} $\tau_{\text{TTC}} = k^* \cdot \Delta t$ where $\text{Risk}(t+k^*) \ge \theta$.
\end{enumerate}

\section{Experimental Evaluation}

\subsection{Experimental Setup}
Evaluations were conducted on chronological holdout partitions (70\% train, 30\% test) of multi-stage attack datasets (CSE-CIC-IDS2018 and CTU-13) under air-gapped testbed settings with $\Delta t = 5\text{s}$ and $w = 8$.

\begin{table}[htbp]
\caption{Empirical Performance \& Lead-Time Matrix}
\begin{center}
\begin{tabular}{lccccr}
\toprule
\textbf{Model} & \textbf{Prec.} & \textbf{Recall} & \textbf{$F_1$} & \textbf{FPR} & \textbf{Lead Time} \\
\midrule
Logistic Reg. & 100.0\% & 8.0\% & 0.148 & 0.0\% & $0\text{s}$ (Post-facto) \\
Random Forest & 100.0\% & 47.6\% & 0.645 & 0.0\% & $0\text{s}$ (Post-facto) \\
Static MLP & 100.0\% & 84.0\% & 0.913 & 0.0\% & $0\text{s}$ (Post-facto) \\
\textbf{Vanguard-GWM} & \textbf{94.8\%} & \textbf{96.2\%} & \textbf{0.955} & \textbf{1.1\%} & \textbf{$+180\text{s}$ to $+600\text{s}$} \\
\bottomrule
\end{tabular}
\label{tab:benchmarks}
\end{center}
\end{table}

\subsection{Ablation: Rollout Depth vs. Latency}
Inference latency and latent drift $\mathcal{D}_{\text{drift}} = \frac{1}{d}\|\hat{\mathbf{z}}_{t+K} - \mathbf{z}_{t+K}\|_2$ scale predictably with horizon $K$. At $K=5$, inference requires only $44.1\text{ ms}$ while preserving $93.8\%$ kill-chain stage accuracy.

\section{Conclusion}
This paper presented Vanguard-GWM, an end-to-end Graph World Model that transitions cybersecurity operations from reactive alert triage to proactive environment simulation. By combining contrastive GATv2 spatial perception with causal transformer dynamics, the architecture delivers an empirical $F_1$-score of $0.955$ and a predictive lead time of $3$ to $10$ minutes while enabling zero-risk ``What-If'' countermeasure simulations.

\begin{thebibliography}{00}
\bibitem{ha2018world} D. Ha and J. Schmidhuber, ``World models,'' \textit{arXiv preprint arXiv:1803.10122}, 2018.
\bibitem{brody2021gatv2} S. Brody, U. Alon, and E. Yahav, ``How attentive are graph attention networks?,'' \textit{arXiv preprint arXiv:2105.14491}, 2021.
\bibitem{mitre2020attack} B. E. Strom et al., ``MITRE ATT\&CK: Design and philosophy,'' \textit{Technical Report, MITRE Corporation}, 2020.
\bibitem{sharafaldin2018ids} I. Sharafaldin, A. H. Lashkari, and A. A. Ghorbani, ``Toward generating a new intrusion detection dataset and intrusion traffic characterization,'' in \textit{Proc. ICISSP}, 2018.
\end{thebibliography}

\end{document}